\chapter{Estudio previo y gestión del trabajo}\label{cap:gestion}

\section{Introducción}
Este capítulo describe cómo se planificó y gobernó el trabajo para llegar a un
resultado técnico completo sin perder trazabilidad. La gestión se orientó a un
objetivo claro: construir una solución reproducible en entorno local, con notebooks
versionados, decisiones justificadas y una narrativa defendible en presentación.

\section{Contexto tecnológico de partida}
La selección tecnológica se apoyó en herramientas abiertas y ampliamente adoptadas
en analítica de datos:
\begin{itemize}
	\item Apache Spark para procesamiento distribuido y consultas analíticas.
	\item Apache Zeppelin para desarrollo exploratorio, análisis y visualización.
	\item Docker Compose para estandarizar el entorno entre miembros del equipo.
	\item Python/PySpark para enriquecimiento geográfico y generación de mapas
	con Folium.
\end{itemize}

La elección de Parquet como formato de entrada permitió mantener coherencia entre
eficiencia de almacenamiento y rendimiento de lectura en consultas analíticas.
Los indicadores por zona se exportan a CSV para su consumo en visualizaciones
geográficas.

\section{Metodología de trabajo}
Se aplicó un enfoque incremental con entregables verificables en cada iteración.
Cada bloque cerraba con una revisión técnica de tres aspectos: viabilidad,
reproducibilidad y claridad documental.

\begin{enumerate}
	\item \textbf{Iteración 1 - Base de ejecución}: estructura de repositorio,
	dockerización y arranque de servicios.
	\item \textbf{Iteración 2 - Carga y limpieza}: notebook principal con carga
	de Parquet, reglas de calidad y enriquecimiento.
	\item \textbf{Iteración 3 - Análisis}: consultas temporales, geográficas y
	económicas en el notebook principal.
	\item \textbf{Iteración 4 - Visualización geográfica}: exportación de CSVs
	por zona y notebook de mapas con Folium.
	\item \textbf{Iteración 5 - Cierre académico}: pruebas, documentación y
	conclusiones.
\end{enumerate}

Este enfoque permitió controlar el alcance y detectar pronto bloqueos técnicos,
evitando una integración tardía de componentes.

\section{Planificación funcional}
La planificación final se alineó con los dos notebooks de Zeppelin y los
artefactos de datos:

\begin{table}[H]
	\centering
	\caption{Planificación por bloques funcionales}
	\begin{tabular}{|p{2.3cm}|p{4.7cm}|p{6.4cm}|}
		\hline
		\textbf{Bloque} & \textbf{Objetivo} & \textbf{Entregable principal} \\
		\hline
		B1 & Ingesta y validación & Carga de Parquet y validación de esquema \\
		\hline
		B2 & Calidad del dato & Reglas de limpieza + dataset limpio en memoria \\
		\hline
		B3 & Análisis temporal & Consultas por hora, día y franja horaria \\
		\hline
		B4 & Análisis geográfico-económico & Join con GeoJSON + indicadores por zona \\
		\hline
		B5 & Exportación de indicadores & CSVs por zona de origen y destino \\
		\hline
		B6 & Mapas interactivos & Notebook de mapas coropléticos con Folium \\
		\hline
	\end{tabular}
\end{table}

\section{Gestión de riesgos}
Durante el trabajo se identificaron riesgos recurrentes en proyectos Big Data
académicos y se definieron respuestas preventivas.

\begin{table}[H]
	\centering
	\caption{Riesgos relevantes y medidas de mitigación}
	\begin{tabular}{|p{4.0cm}|p{3.2cm}|p{6.2cm}|}
		\hline
		\textbf{Riesgo} & \textbf{Impacto} & \textbf{Mitigación aplicada} \\
		\hline
		Memoria insuficiente en local & Alto & Análisis centrado en un mes (enero 2025) y ajuste de parámetros Spark \\
		\hline
		Inconsistencia entre entornos del equipo & Alto & Uso de Docker Compose y comandos estandarizados \\
		\hline
		Resultados difíciles de defender & Medio & Reglas de limpieza explícitas y trazabilidad de salidas \\
		\hline
		Dependencia de datos auxiliares geográficos & Medio & Generación de GeoJSON desde shapefile con script auxiliar \\
		\hline
		Sobrecarga documental al final & Medio & Redacción continua por iteraciones y cierre por capítulo \\
		\hline
	\end{tabular}
\end{table}

\section{Recursos y estimación de esfuerzo}
El proyecto se desarrolló en contexto docente, por lo que no se imputan licencias
de software. La estimación económica se centra en esfuerzo de equipo y coste operativo
del entorno local.

\begin{table}[H]
	\centering
	\caption{Estimación orientativa de esfuerzo y coste}
	\begin{tabular}{|p{6.8cm}|c|c|}
		\hline
		\textbf{Concepto} & \textbf{Cantidad} & \textbf{Coste estimado (EUR)} \\
		\hline
		Esfuerzo técnico del equipo (análisis, código y documentación) & 120 h & 0 \\
		\hline
		Software y licencias (Spark, Zeppelin, Python, Docker) & -- & 0 \\
		\hline
		Infraestructura cloud & -- & 0 \\
		\hline
		Consumo eléctrico y uso de equipos locales & -- & 20 \\
		\hline
		\textbf{Total estimado} & -- & \textbf{20} \\
		\hline
	\end{tabular}
\end{table}

\section{Conclusiones}
La gestión del trabajo permitió mantener equilibrio entre ambición técnica y alcance
realista. El resultado no es un prototipo aislado, sino una solución completa con
notebooks de análisis y mapas, exportación de indicadores y documentación coherente,
preparada para presentación académica y para evolución futura a entornos de mayor
escala.